\clearpage
\begin{explanationbox}
\textbf{先把双曲线拿掉：熟悉的角平分线定理。}
\begin{center}
\begin{tikzpicture}[font=\small,>=stealth]
\draw[gray!40,rounded corners] (0,0) rectangle (14.2,8);
\node[font=\small\bfseries] at (7.1,7.35) {第一步：只有普通三角形，不需要双曲线};

\begin{scope}[shift={(4.5,3.1)},scale=.62]
\coordinate (A) at (-5,0); \coordinate (B) at (5,0);
\coordinate (C) at (3.75,3); \coordinate (D) at (2.4,0);
\fill[blue!8] (A)--(C)--(D)--cycle;
\fill[orange!12] (B)--(C)--(D)--cycle;
\draw[blue!75!black,very thick] (A)--(C) node[midway,above left] {$u$};
\draw[orange!85!black,very thick] (C)--(B) node[midway,right] {$v$};
\draw[blue!75!black,very thick] (A)--(D) node[midway,below] {$AD$};
\draw[orange!85!black,very thick] (D)--(B) node[midway,below] {$DB$};
\draw[magenta!80!black,very thick] (C)--(D);
\draw (C) ++(-161.075:1) arc (-161.075:-114.228:1);
\draw (C) ++(-114.228:1.2) arc (-114.228:-67.38:1.2);
\foreach \q in {A,B,C,D} \fill (\q) circle(3pt);
\node[below left] at (A) {$A$};\node[below right] at (B) {$B$};
\node[above] at (C) {$C$};\node[below,yshift=-20pt] at (D) {$D$};
\end{scope}
\node[align=left,anchor=west] at (9,5.3) {$\angle ACD=\angle DCB$\\[10pt]$\displaystyle\frac{AD}{DB}=\frac{u}{v}$};
\node[align=center] at (7.1,1) {同色对应：左边对左段，右边对右段。\\下一步只需换名：$A,B,C,D\ \longrightarrow\ F_1,F_2,P,M$。};

\end{tikzpicture}
\end{center}


\textbf{再把同一结构放回焦点三角形。}
将图中 $A,B,C,D$ 换成 $F_1,F_2,P,M$，令 $PF_1=r_1,PF_2=r_2$。
由于内角平分线交点在底边内部，$F_1M=c+m,MF_2=c-m$，于是
\[
\frac{c+m}{c-m}=\frac{r_1}{r_2}
\quad\Longrightarrow\quad
\boxed{m=c\frac{r_1-r_2}{r_1+r_2}}.
\]
这一步只用初中几何，尚未使用双曲线性质。

\textbf{最后才加入双曲线的两条信息：}
\[
r_1^2-r_2^2=4cx_0,\qquad |r_1-r_2|=2a.
\]
前式来自勾股定理与平方差（见方法一配图），后式来自双曲线定义。
把它们与角平分线所得的式子相乘，即得
\[
\boxed{mx_0=\frac{(r_1-r_2)^2}{4}=a^2}.
\]
先求出交点 $M$，再由 $P,M$ 两点写出直线，便能独立证明它就是切线。
内切圆切点 $T$ 与交点 $M$ 的区别，见方法四配图。

\medskip
\textbf{另一条独立路线：倒角公式（正切差角公式）。}
若 $k_1,k_2,k$ 分别为 $PF_1,PF_2,PM$ 的斜率，
内角平分线给出相等的两个半角。在斜率存在的普通情形下，
\[
\boxed{\frac{k-k_1}{1+kk_1}=\frac{k_2-k}{1+k_2k}},\qquad
k_1=\frac{y_0}{x_0+c},\quad
k_2=\frac{y_0}{x_0-c},\quad
k=\frac{y_0}{x_0-m}.
\]
代入并利用 $y_0^2=\dfrac{b^2}{a^2}(x_0^2-a^2)$，因式分解得
\[
(mx_0-a^2)(c^2x_0-a^2m)=0.
\]
两个候选交点分别为 $m=a^2/x_0$（\textbf{内角平分线}）和
$m=c^2x_0/a^2$（\textbf{外角平分线}）。后者满足 $|m|>c$，
不在两个焦点之间，因此不能当作题设中的 $M$。
\textbf{遇到竖直焦点连线时}不能硬套斜率，而要改用向量点积与行列式，
细节见证明方法六。
\end{explanationbox}
